\scanpage{226}
matrices \unsure{$\{$} $\underline u$ \unsure{$\in B$} telles que $u\equiv1\ (2)$ i.e. $M\equiv0\ \unsure{(2)}\unsure{\}}$, par le ss-groupe \inserted{fermé} $\underline L_0$ engendré par $\ell_0=\begin{pmatrix}1&-2\\0&1\end{pmatrix}$, s'identifie au groupe des matrices $\begin{pmatrix}\mu&0\\0&1\end{pmatrix}$ --- \inserted{\unsure{le}} tore maximal $T_0$ de $B$, de façon plus précise on a une structure de produit $\frac12$ direct
\[
\begin{array}{ccccc}
\unsure{\underline B_1}&\simeq&T_0&\cdot&\underline L_0\\
\Vert&&\Vert&&\Vert\\
\left\{\begin{pmatrix}\mu&M\\0&1\end{pmatrix}\middle|\begin{array}{l}\mu\in\widehat{\mathbb Z}^{*}\\M\in2\widehat{\mathbb Z}\end{array}\right\}&&
\left\{\begin{pmatrix}\mu&0\\0&1\end{pmatrix}\middle|\mu\in\widehat{\mathbb Z}^{*}\right\}&&
\left\{\begin{array}{c}\ell_0^n\\\Vert\\\begin{pmatrix}1&-2n\\0&1\end{pmatrix}\end{array}\middle|n\in\widehat{\mathbb Z}\right\}
\end{array}
\tag{178}
\]

Cela montre que si on s'intéresse \unsure{au quotient|aux quotients} $\mathcal M_{0,3}^{!\sim}\unsure{(0)|[0]}/L_0\xrightarrow{\sim}\mathcal N_{0,3}^{!\sim}$, alors le seul invariant abélien qu'on sache lui associer dans le contexte présent (quand on ne travaille pas avec $\alpha,\beta$) est le multiplicateur $\mu$. Mais il y a mieux\unsure{${}^{\times}$}, on trouve \unsure{grâce} à la décomposition (178) dans $\mathrm{GL}(2,\widehat{\mathbb Z})$ un scindage canonique de l'extension $\mathcal M_{0,3}^{!\sim}(0)$ de $\mathcal N_{0,3}^{!\sim}$ par $L_0$ :

\uline{Proposition} Pour $\forall\,\dot u\in\mathcal N_{0,3}^{!\sim}$, il existe un représentant unique \unsure{$u=u_{\dot u}$|$u=u_{\xi}$} $\in\mathcal M_{0,3}^{!\sim}[0]$ tel qu'on ait $\underline u\in T_0$ i.e. $\underline u$ de la forme $\begin{pmatrix}\mu&0\\0&1\end{pmatrix}$
\inserted{[i.e. $\xi=\begin{pmatrix}A&B\\C&D\end{pmatrix}$ de la forme $\begin{pmatrix}A&0\\C&D\end{pmatrix}=\begin{pmatrix}\unsure{\frac{1}{D}}&0\\1-D&D\end{pmatrix}$]}, --- En d'autres termes, le sous-groupe

\margin{225}

\scanpage{227}
\mspage{413}
\deleted{\unsure{$\mathcal M_{0,3}^{!\sim}(0)$}}
\[
\mathcal M_{0,3}^{!\sim}(0)\overset{\mathrm{df}}=
\left\{u\in\mathcal M_{0,3}^{!\sim}\unsure{(0)|[0]}\ \middle|\ \underline u\in T_0\quad\text{i.e. }\underline u=\begin{pmatrix}\mu&0\\0&1\end{pmatrix}\right\}
\tag{179}
\]
est un ss-groupe section de l'extension $\mathcal M_{0,3}^{!\sim}[0]$ de $\mathcal N_{0,3}^{!\sim}$ par $L_0$, donc aussi de l'extension $\mathcal M_{0,3}^{!\sim}$ de $\mathcal N_{0,3}^{!\sim}$ par $\widehat\Pi_{0,3}$, \unsure{ou} de l'extension $\mathcal M_{0,3}^{\sim}$ de $\mathcal N_{0,3}^{!\sim}$ par $\widehat{\mathcal G}_{0,3}^{+}$.

\unsure{Ceci|Cela} provient du fait que si $u\in\mathcal M_{0,3}^{!\sim}[0]$ \inserted{définit} $\underline u=\begin{pmatrix}\mu&M\\0&1\end{pmatrix}$, avec donc $M\equiv0\ (2)$, i.e. $M\in2\widehat{\mathbb Z}$, il existe un \uline{unique} $n\in\widehat{\mathbb Z}$ tel que $\underline{\ell_0^n u}=\begin{pmatrix}\mu&M-2n\\0&1\end{pmatrix}$ soit dans $T_0$ i.e. tel que $\ell_0^n u\in\mathcal M_{0,3}^{!\sim}(0)$, savoir $n=\dfrac M2$.

Si $(\xi,\deleted{\unsure{\varepsilon}}\eta)$ satisfont aux conditions (133)(134) i.e. définissent un $u=u_{\xi,\eta}$ (cf Th. p.~407'), alors on dira que $(\xi,\eta)$ est \uline{normalisé} si on a $u_{\xi,\eta}\in\mathcal M_{0,3}^{!\sim}(0)$. De $\overline{\mathrm m}$, si un élément quelconque de $\mathcal M_{0,3}^{!\sim}$ est écrit sous la forme
\[
u=\cdot f^{-1}\cdot u_{\xi,\eta}=f^{-1}u_{\unsure{\xi}},
\]
on dit que cette représentation \inserted{(\unsure{entachée} de l'indétermination (136))} est \uline{normalisée}, si $u_{\xi,\eta}=u_\xi$ est normalisé, [i.e. \unsure{si} \inserted{pour} $\xi\unsure{\simeq}\begin{pmatrix}A&B\\C&D\end{pmatrix}$ \deleted{la forme} on a $B=0$]. Ainsi, \deleted{tout} $u$ se met de façon unique sous forme
\[
u_{f,\xi,\eta}=f^{-1}u_\xi,
\tag{\unsure{179}}
\]
\margin{226}

\scanpage{228}
avec $\unsure{f|\xi}$ normalisé. En fait, ceci correspond à la décomposition en produit $\frac12$-direct
\[
\mathcal M_{0,3}^{!\sim}\simeq\mathcal M_{0,3}^{!\sim}(0)\cdot\widehat\Pi_{0,3},
\tag{180}
\]
on a d'ailleurs de $\overline{\mathrm m}$
\[
\mathcal M_{0,3}^{\sim}\simeq\mathcal M_{0,3}^{!\sim}(0)\cdot\widehat{\mathcal G}_{0,3}^{+}.
\tag{181}
\]

\uline{Rmq} Il faudrait à l'occasion \unsure{essayer} de mettre ces scindages canoniques en relation avec le choix d'une uniformisante en $0$ de $\mathbb P^1_{\mathbb Q}=\widehat U_{0,3\mathbb Q}$. Bien entendu, la $\overline{\mathrm m}$ question se pose aussi \unsure{sur|pour} les \inserted{autres deux} ``points à l'infini'' $1,\infty$ de $U_{0,3\mathbb Q}$, relativement aux ss-groupes sections
\[
\mathcal M_{0,3}^{!\sim}(q)\qquad q\in\{0,1,\infty\}
\tag{182}
\]
correspondant à ces points à l'infini --- la construction d'un groupe-section $\mathcal M_{0,3}^{!\sim}(0)$ dans le normalisateur de $L_0$, s'applique évidemment sans chgt pour $L_1,L_\infty$. On aura d'ailleurs
\[
\rho\bigl(\mathcal M_{0,3}^{!\sim}(q)\bigr)=\mathcal M_{0,3}^{!\sim}\bigl(\unsure{\dot\rho|\rho}(q)\bigr)\qquad\text{pour }q\in\{0,1,\infty\},
\tag{183}
\]
donc les \inserted{trois} $\mathcal M_{0,3}^{!\sim}(q)$ se déduisent de l'un quelconque d'entre eux en appliquant $\rho,\rho^2$.
\margin{227}

\scanpage{229}
\mspage{415}
47. Retour sur l'homomorphisme $\widehat\Pi_{0,3}\xrightarrow{\widehat\varphi}\widehat{\mathcal G}_{0,3}^{+}$ et la condition de commutation extérieure à \uline{celui-ci}. Les groupes $\mathcal M_{0,3}^{!\approx},\mathcal M_{0,3}^{\approx}$ et l'hom. $\Phi:\mathcal M_{0,3}^{!\approx}\longrightarrow\mathcal M_{0,3}^{\approx}$.

Rappelons \inserted{(\S41)} que $\widehat\varphi$ est décrit par
\[
\widehat\varphi(\ell_0)=\varepsilon_0=\sigma_\infty\rho,\quad
\widehat\varphi(\ell_1)=\sigma_\infty,\quad
\widehat\varphi(\ell_\infty)=\rho^{-1}.
\tag{1}
\]
Soit
\[
u=u_{\alpha,\beta,f}\quad\in\mathcal M_{0,3}^{!\sim}
\]
un autom. à lacets de $\widehat\Pi_{0,3}$, commutant ext${}^{t}$ à $\mathcal G_3$, défini par
\[
\begin{aligned}
u(\sigma_\infty)&=h\sigma_\infty=\operatorname{int}(\beta_0)\sigma_\infty
&\left(\begin{array}{l}h=\beta_0\sigma_\infty(\beta_0)^{-1},\\\beta_0=f^{-1}\beta\end{array}\right)
\end{aligned}
\tag{2}
\]
\[
\begin{aligned}
u(\rho)&=g\rho=\operatorname{int}(\alpha_0)\rho
&\left(\begin{array}{l}g=\alpha_0\rho(\alpha_0)^{-1},\ \deleted{\unsure{=}}\\\alpha_0=f^{-1}\alpha\end{array}\right)
\end{aligned}
\tag{3}
\]
et écrivons que $u$ commute extérieurement à $\varphi$, i.e. qu'il existe \deleted{\unsure{$\lambda\in$}}
\[
\lambda\in\widehat\Pi_{0,3}
\tag{4}
\]
tel que l'on ait
\[
\widehat\varphi u(x)=\operatorname{int}(\lambda)u\widehat\varphi(x)
\qquad\forall x\in\widehat\Pi_{0,3}.
\tag{5}
\]
Il suffit de vérifier (5) pour $x=\ell_0$ et $x=\ell_1$.

1) Cas $x=\ell_0$ \deleted{\unsure{$=(\sigma_\infty\rho)^2$}}, d'où
\[
u(x)=\deleted{u(\sigma_\infty)u(\rho)}\deleted{\unsure{=h\sigma_\infty g\rho=h(\sigma_\infty g\sigma_\infty^{-1})\sigma_\infty\rho}}
\inserted{\operatorname{int}(f^{-1})(\ell_0^\mu)}
\]
\[
\widehat\varphi(u(x))=\operatorname{int}\bigl(\widehat\varphi(f)^{-1}\bigr)\widehat\varphi(\ell_0^\mu)=\operatorname{int}\bigl(\widehat\varphi(f)^{-1}\bigr)\varepsilon_0^\mu
\]
\[
u\widehat\varphi(x)=u(\varepsilon_0)=\operatorname{int}(f^{-1})(\varepsilon_0^\mu)
\]
donc la relation (5) s'écrit
\[
\operatorname{int}\bigl(\deleted{\unsure{f\lambda^{-1}}}\widehat\varphi(f)\unsure{\lambda}f^{-1}\bigr)(\varepsilon_0^\mu)=\varepsilon_0^\mu
\]
\margin{228}

\scanpage{230}
équivaut à
\[
\deleted{\unsure{f\lambda^{-1}}}\widehat\varphi(f)\lambda f^{-1}\in L_0^{\sim}\quad
\left(=\operatorname{Centr}_{\unsure{\widehat\Pi_{0,3}}}(\varepsilon_0^\mu)\right)
\]
\margin{NB on n'a pas nécessairement $\widehat\varphi(f)\in\widehat\Pi_{0,3}$, donc \unsure{pas non plus} \illegible{1} $\widehat\varphi(f)\lambda f^{-1}\in\widehat\Pi_{0,3}\ldots$}
i.e. à l'existence de $r\in\widehat{\mathbb Z}$ tel que l'on ait
\[
\deleted{\unsure{f\lambda^{-1}}}\widehat\varphi(f)\lambda f^{-1}=\varepsilon_0^r,
\]
i.e.
\[
\lambda=\widehat\varphi(f)^{-1}\varepsilon_0^r f.
\tag{6}
\]
Ceci permet donc d'éliminer $\unsure{\exists\lambda}$ de la deuxième relation :

2) Cas $x=\ell_1$, d'où
\[
u(x)=\operatorname{int}(f_1^{-1})(\ell_1^\mu)
\qquad\bigl(f_1^{-1}=\deleted{\unsure{\rho(f)}}g\rho(f^{-1})\bigr)
\]
\[
\widehat\varphi(u(x))=\operatorname{int}\bigl(\widehat\varphi(f_1)^{-1}\bigr)\sigma_\infty^\mu=\operatorname{int}\bigl(\widehat\varphi(f_1)\bigr)\sigma_\infty
\]
(compte tenu que $\sigma_\infty^\mu=\sigma_\infty$ car $\mu\equiv1\ (2)$).
\deleted{\unsure{$u\widehat\varphi(x)=\operatorname{int}(\widehat\varphi(f_1))$}}
\[
u\widehat\varphi(x)=u(\sigma_\infty)=\operatorname{int}(\beta_0)(\sigma_\infty),
\]
Donc la relation (5) s'écrit
\[
\operatorname{int}\bigl(\widehat\varphi(f_1)^{\unsure{+1|-1}}\lambda\beta_0\bigr)\sigma_\infty=\sigma_\infty,
\]
ce qui signifie aussi qu'il existe $\unsure{i|\alpha}\in\mathbb Z/2\mathbb Z$ tel qu'on ait
\[
\widehat\varphi(f_1)^{\unsure{+1|-1}}\lambda\beta_0=\sigma_\infty^i
\quad\text{soit}\quad
\widehat\varphi(f_1^{-1})^{\unsure{+1|-1}}=\lambda\beta_0\sigma_\infty^i,
\]
et remplaçant $\lambda$ par sa valeur (6), et $f_1^{-1}$ par sa valeur $g\rho(f^{-1})\deleted{\unsure{f}}$, on trouve la condition
\[
\deleted{\unsure{\widehat\varphi(\rho(f)^{-1})\widehat\varphi(g)^{-1}=\widehat\varphi(f)^{-1}\ell_0^r f\beta_0\sigma_\infty^\alpha}}
\]
\[
\widehat\varphi\bigl(g\rho(f)^{-1}\bigr)=\widehat\varphi(f)^{-1}\varepsilon_0^r\underbrace{f\beta_0}_{\beta}\sigma_\infty^{\inserted{\unsure{\alpha}}}\quad\text{soit}
\]
\[
\widehat\varphi\bigl(\underbrace{fg\rho(f)^{-1}}_{\xi=\alpha\rho(\alpha)^{-1}}\bigr)=\varepsilon_0^r\beta\sigma_\infty^\alpha
\]
\margin{229}

\scanpage{231}
\mspage{416}

Ainsi :

\uline{\deleted{Proposition}\inserted{Théorème}} Pour que
\[
\widehat\varphi:\widehat\Pi_{0,3}\longrightarrow
\widehat{\unsure{\mathcal G}}_{0,3}^{+}
\]
commute extérieurement à l'élément
$u\in\mathcal M_{0,3}^{!\sim}$, correspondant aux invariants
$g,h,f,\alpha_0,\beta_0,\alpha,\beta,\xi,\eta$ (cf \S 46),
il f. et s. qu'on ait
\[
\widehat\varphi(\xi)=\varepsilon_0^r\beta\sigma_\infty^{\unsure{i|\alpha}}
\qquad\text{pour }r\in\widehat{\mathbb Z},\quad
\alpha\in\mathbb Z/2\mathbb Z\ \text{convenables}.
\tag{7}
\]
\margin{\uline{NB} $r,\alpha$ sont bien déterminés quand $\alpha$ et $\beta$
(ou \unsure{encore} $\xi$ et $\beta$) sont choisis.}
On aura alors (5), avec \deleted{Notons que}
$\lambda$ donné par (6) ;
$\lambda=\widehat\varphi(f)^{-1}\ell_0^r f$.

Notons que $u$ ne définit le triple $(\alpha,\beta,f)$ que modulo
une transformation de la forme
\deleted{$(\alpha,\beta,f)\ \unsure{\mapsto}\ \unsure{\alpha'=\ell_0^n\alpha,\ \beta}$}
\[
(\alpha,\beta,f)\longmapsto
(\alpha'=\ell_0^n\alpha,\ \beta'=\ell_0^n\beta,\ f'=\ell_0^n f)
\tag{8}
\]
où $n\in\widehat{\mathbb Z}$. On aura alors (posant
$\xi'=\alpha'\rho(\alpha')^{-1}=\ell_0^n\xi\ell_1^{-n}$)
\[
\widehat\varphi(\xi')=\varepsilon_0^{r'}\beta'
\sigma_\infty^{\unsure{i'|\alpha'}}
\tag{*}
\]
avec $r'\in\widehat{\mathbb Z}$,
$\unsure{i'|\alpha'}\in\mathbb Z/2\mathbb Z$ également bien déterminés.
Or on aura
\[
\begin{aligned}
\widehat\varphi(\unsure{\xi'})
 &=\widehat\varphi(\ell_0^n\unsure{\xi}\ell_1^{-n})\\
 &=\widehat\varphi(\ell_0^n)\widehat\varphi(\unsure{\xi})
   \widehat\varphi(\ell_1^{-n})\\
 &=\varepsilon_0^n\widehat\varphi(\unsure{\xi})\sigma_\infty^{-n}
\end{aligned}
\]
\deleted{\illegible{1}}\deleted{\illegible{1}}
et la relation (*) s'écrit
\[
\widehat\varphi(\xi)=\varepsilon_0^{\unsure{r'-n|r'-2n}}\beta'
\sigma_\infty^{\unsure{i'-n|\alpha'-n}}
\]
ou encore, comme $\beta'=\ell_0^n\beta=\varepsilon_0^{2n}\beta$
\[
\widehat\varphi(\xi)=\varepsilon_0^{r'+n}\beta
\sigma_\infty^{\unsure{\alpha'-n|i'-n}},
\]
ce qui, comparé à (7), donne
\[
r'=r-n,\qquad\unsure{i'|\alpha'}=\unsure{i|\alpha}+n\ \mathrm{mod}\,2.
\tag{9}
\]
\margin{230}

\scanpage{232}

Cela implique le

\uline{Corollaire 1} Supposons que $u$ commute extérieurement à
$\widehat\varphi:\widehat\Pi_{0,3}\longrightarrow
\widehat{\mathcal G}_{0,3}^{+}$. Alors il est possible
\uline{de façon unique} de décrire \unsure{$u$|$\varphi$} par un triple
$(\alpha,\beta,f)$ (a priori déterminé mod une transformation du type (8),
de façon à ce qu'on ait (posant $\xi=\alpha\rho(\alpha)^{-1}$,
comme à l'accoutumé) \unsure{une} relation
\[
\boxed{\widehat\varphi(\xi)=\beta\sigma_\infty^i}
\qquad\text{pour }i\in\mathbb Z/2\mathbb Z\ \text{convenable}.
\tag{10}
\]
\margin{Il faudrait examiner si cette définition coïncide avec celle de
\unsure{la} p.\,\unsure{410'}\,,\,\unsure{413|417}, formulée sans faire
intervenir $\widehat\varphi$.}
Un tel choix \unsure{sera} appelé \uline{normalisé}.

Ainsi les $u\in\mathcal M_{0,3}^{!\sim}$ qui commutent
extérieurement à $\widehat\varphi$ correspondent biunivoquement aux
\deleted{couples} triples \deleted{$(\alpha,\beta,f)$}
\deleted{satisfaisant la (10)}
\[
(\alpha,\beta,f),\qquad
\alpha\in\unsure{\widetilde{\widehat\Pi}_{0,3}},\quad
\beta,f\in\widehat\Pi_{0,3}
\]
satisfaisant, d'une part la relation des bouts
\[
\xi=\eta\ell_1^\nu
\qquad
\begin{gathered}
\text{où }\xi=\alpha\rho(\alpha)^{-1},\quad
\eta=\beta\sigma_\infty(\beta)^{-1},\\
\text{et où }\nu\in\widehat{\mathbb Z}\ \text{convenable}\\
(\text{tel que }2\nu+1\in\widehat{\mathbb Z}^{*})
\end{gathered}
\tag{11}
\]
d'autre part la relation (10) (pour $i\in\mathbb Z/2\mathbb Z$ conv.)
--- l'une et l'autre relation ne faisant d'ailleurs pas intervenir $f$.
\margin{231}

\scanpage{233}
\mspage{417}

Notons que $i\in\mathbb Z/2\mathbb Z$ est défini en termes de
$\xi$ (compte tenu de
$\beta\in\widehat\Pi_{0,3}=\operatorname{Ker}
(\widehat{\mathcal G}_{0,3}^{+}\to\unsure{\mathcal G_3|\mathfrak S_3})$)
en réduisant (10) mod $\widehat\Pi_{0,3}$. Désignant par
\[
\widehat\varphi^{\unsure{\bullet|\circ}}:
\widehat\Pi_{0,3}\longrightarrow\unsure{\mathcal G_3|\mathfrak S_3}
\tag{12}
\]
le composé
\[
\widehat\Pi_{0,3}\xrightarrow{\widehat\varphi}
\widehat{\mathcal G}_{0,3}^{+}
\overset{\mathrm{can}}{\unsure{\twoheadrightarrow|\longrightarrow}}
\unsure{\mathcal G_3|\mathfrak S_3},
\]
la formule (10) donne aussi
\[
\widehat\varphi^{\unsure{\bullet|\circ}}(\xi)=\sigma_\infty^i
\tag{13}
\]
ce qui détermine bien $i$\inserted{$\in\mathbb Z/2\mathbb Z$} en termes de
$\xi\in\widehat\Pi_{0,3}$, et implique une condition intrinsèque sur
l'élément $\xi$ de $\widehat\Pi_{0,3}$, à savoir que son image dans
$\unsure{\mathcal G_3|\mathfrak S_3}$ par
$\widehat\varphi^{\unsure{\bullet|\circ}}$ soit bien dans le sous-groupe
$\{1,\sigma_\infty\}$ de $\unsure{\mathcal G_3|\mathfrak S_3}$.
En fait, le noyau de $\widehat\varphi^{\unsure{\bullet|\circ}}$ est
le sous-groupe \unsure{inv.} de $\widehat\Pi_{0,3}$ engendré par
$\ell_0^2,\ell_1^2,\ell_\infty^3$, et la condition envisagée signifie
donc que modulo ce sous-groupe, $\xi$ soit de la forme $\ell_1^i$.

Par la formule (10), on peut éliminer $\beta$ de la donnée
$(\alpha,\beta,f)$, et on peut remplacer la donnée $\alpha$ par la donnée
$\xi$. Ainsi :
\margin{232}

\scanpage{234}

\uline{Corollaire 2} \unsure{Les} éléments du \unsure{$^{s}/$groupe} de
$\mathcal M_{0,3}^{!\sim}$ formé des $u$ qui commutent extérieurement à
$\widehat\varphi$, est \deleted{\illegible{1}} en corr. biunivoque avec
l'ens des couples $(\xi,f)$ \deleted{\unsure{$\in$}}
$(\xi,f\in\widehat\Pi_{0,3})$ satisfaisant les conditions
(\unsure{ne|en} faisant intervenir que $\xi$)
\[
\begin{array}{c@{\qquad}c}
\begin{gathered}(14)\\[1.8ex](15)\end{gathered}
&
\boxed{\begin{gathered}
\xi\rho(\xi)\rho^2(\xi)=1\\[1.8ex]
\xi=\bigl(\widehat\varphi(\xi)\,
\sigma_\infty(\widehat\varphi(\xi))^{-1}\bigr)\ell_1^\nu
\end{gathered}}
\end{array}
\]
\[
(\nu\in\widehat{\mathbb Z}\ \text{convenable tel que }
2\nu+1\in\widehat{\mathbb Z}^{*}).
\]
\margin{il faut de plus la ``condition'' ``\unsure{liminaire}''\ldots}

Pour un tel $\xi$, $\widehat\varphi^{\unsure{\bullet|\circ}}(\xi)$
dans $\unsure{\mathcal G_3|\mathfrak S_3}$ est de la forme
$\sigma_\infty^i$, avec $i\in\mathbb Z/2\mathbb Z$, et posant
\[
\left\{\begin{aligned}
&\beta=\widehat\varphi(\xi)\sigma_\infty^i
  \quad(\in\widehat\Pi_{0,3})\\
&\alpha\in\widehat\Pi_{0,3}\quad
  \text{l'unique élément de }\widetilde{\widehat\Pi}_{0,3}
  =\widehat\Pi_{0,3}\cdot\{1,\tau_\infty\}\\
&\hspace{5em}\text{tel que }\alpha\rho(\alpha)^{-1}=\xi
\end{aligned}\right.
\tag{16}
\]
\margin{en fait, on n'a pas vraiment besoin de $\alpha,\beta$}
$u$ se déduit de $\xi,f$ par les formules
\[
\left\{\begin{aligned}
u(\sigma_\infty)
 &=\operatorname{int}(f^{-1}\unsure{\beta})(\sigma_\infty)\\
 &=\operatorname{int}(f^{-1})(\operatorname{int}(\unsure{\beta})(\sigma_\infty))\\
 &=\operatorname{int}(f^{-1})(\xi\ell_1^{-\nu}\sigma_\infty)\\
u(\rho)
 &=\operatorname{int}(f^{-1}\alpha)(\rho)\\
 &=\operatorname{int}(f^{-1})(\operatorname{int}(\alpha)(\rho))\\
 &=\operatorname{int}(f^{-1})(\unsure{\xi}\rho)
\end{aligned}\right.
\tag{17}
\]

On notera que dans ces formules, $f$ n'intervient que par le fait de
composer avec $\operatorname{int}(f)$. Si on
\deleted{pose}\inserted{désigne}, pour $\xi\in\widehat\Pi_{0,3}$
\margin{233}

\scanpage{235}
\mspage{418}

\margin{Cette \unsure{notation} est déjà de mise pour les
\unsure{considérations} générales du \S précédent
$\unsure{u_{\xi,\eta}}$ ne dépend de \unsure{toute} façon que de $\xi$\ldots}
satisfaisant (14)(15), par $u_\xi$ l'autom.
\deleted{$u_\xi(\sigma_\infty)$} correspondant, en termes du
Corollaire 2, au couple $(\xi,1)$, de sorte que
\[
\left\{\begin{aligned}
u_\xi(\sigma_\infty)
 &=\deleted{\unsure{\operatorname{int}(\alpha)}}
 \overbrace{(\xi\ell_1^{-\nu})}^{\unsure{\eta|\nu}}\sigma_\infty\\
u_\xi(\rho)&=\xi\rho
\end{aligned}\right.,
\tag{18}
\]
les formules (17) s'écrivent aussi, simplement
\[
u_{\xi,f}=\operatorname{int}(f^{-1})\,u_\xi.
\tag{19}
\]

D'autre part les formules (5), (6) donnent ici
\[
\widehat\varphi\circ u_\xi=u_\xi\circ\widehat\varphi
\tag{20}
\]
où, bien sûr, dans le premier \unsure{mbr|mb} $u_\xi$ désigne la
restriction de $u_\xi$ à $\widehat\Pi_{0,3}$.

Notons

\uline{Corollaire 3} Les automorphismes à bouts de
$\widehat{\unsure{\mathcal G}}_{0,3}^{+}$ de la forme $u_\xi$
\inserted{($\xi$ \unsure{normalisé|normalisés}$^*$)} sont exactement ceux pour
lesquels on a les \deleted{propriétés} formules (20) (i.e. qui
\deleted{\illegible{1}} commutent à $\widehat\varphi$) et qui de plus
normalisent $L_0$.
\margin{$^*$ au sens de p.\,416'}
\figure{\begin{tikzpicture}[x=1mm,y=1mm,line width=.3pt,line cap=round]
\draw (0,35) .. controls (-.16,29) and (.15,23) .. (-.02,17);
\draw (.65,35.0) -- (.88,34.1) -- (1.38,34.55);
\draw (0,9.5) .. controls (-.08,6.3) and (-.16,2.7) .. (0,0);
\end{tikzpicture}}
\margin{234}

\scanpage{236}

En effet, en termes des \inserted{triples} $(\alpha,\beta,f)$ \inserted{``spéciaux''} ces propriétés s'expriment par les conditions \inserted{\unsure{simples}} \deleted{que l'on ait} $f$ de la forme $\ell_0^n$ (qui exprime la condition de normaliser $L_0$) et que $\lambda$ donné par (6) \uline{\unsure{car|cas}} $r=0$, i.e.
\[
\lambda=\widehat\varphi(f)^{-1}f
\tag{21}
\]
soit égal à $1$, i.e. la condition
\[
\widehat\varphi(f)=f,
\tag{22}
\]
qui en l'occurrence s'écrit
\[
\varepsilon_0^n=\underbrace{\ell_0^n}_{\varepsilon_0^{2n}}
\qquad\text{i.e.}\quad \varepsilon_0^n=1
\]
ce qui implique $\unsure{f|f_0}=\ell_0^n=1$, donc $u$ de la forme $u_{(\alpha,\beta,1)}=u_\xi$, où $\xi=\alpha\rho(\alpha)^{-1}$.

\uline{Corollaire 4} Les $u_\xi$ \inserted{normalisés} forment un \uline{sous-groupe} de $\operatorname{Norm}_{\mathcal M_{0,3}^{!\sim}}(L_0)$, soit $\mathcal M_{0,3}^{!\approx}(0)$, dont l'\deleted{\unsure{image}}\inserted{intersection} avec $\widehat\Pi_{0,3}$ est réduit à $\{1\}$, et dont l'image dans $\mathcal N_{0,3}^{!\sim}$ est \inserted{le ss-groupe $\mathcal N_{0,3}^{!\approx}$} formé des automorphismes ext. à \unsure{la carte|la cat.} de $\widehat{\mathcal G}_{0,3}^{+}$, stabilisant $\widehat\Pi_{0,3}$ (i.e. $u\in\unsure{\mathcal M_{0,3}^{+\widehat{\ }}}$) tels que $\widehat\varphi:\widehat\Pi_{0,3}\to\widehat{\mathcal G}_{0,3}^{+}$ commute extérieurement à $u$. On a donc un isomorphisme
\[
\mathcal M_{0,3}^{!\approx}(\unsure{0})
=\{u_\xi\mid \xi\in\widehat\Pi_{0,3}\text{ solution de }(14)(15)\}
\xrightarrow{\sim}\mathcal N_{0,3}^{!\approx}.
\tag{\unsure{23}}
\]
\margin{235}

\scanpage{237}
\mspage{419}

Il reste seulement à vérifier qu'un $u_\xi$ ne peut être un autom intérieur que si $u_\xi=1$ i.e. $\xi=1$ (donc $\nu=0$). Or $u_\xi$ est un $u_{\alpha,\beta}=u_{\alpha,\beta,1}$ ($\alpha,\beta$ définis \deleted{par} via $\xi$ par $\alpha\rho(\alpha)^{-1}=\xi$, $\beta\sigma_\infty(\beta)^{-1}=\xi \ell_1^{-\nu}$), qui est un autom int. ssi on a $\alpha=\ell_0^n$, $\beta=\ell_0^n$ (cf prop. de la page 404), i.e. $\xi=\ell_0^n \ell_1^{-n}$, $\eta=\beta\sigma_\infty(\beta)^{-1}=\ell_0^n \ell_1^{-n}$, et \deleted{la relation} \deleted{\unsure{pour}} $\xi=\eta \ell_1^\nu$ s'écrit $\ell_1^\nu=1$ i.e. $\nu=0$. La relation (15) devient donc (comme
\[
\begin{gathered}
\widehat\varphi(\xi)=\widehat\varphi(\ell_0^n \ell_1^{-n})=\varepsilon_0^n\sigma_\infty^n\\
\sigma_\infty(\widehat\varphi(\xi))=\varepsilon_1^n\sigma_\infty^n
\end{gathered}
\]
) \deleted{\unsure{donne}} la relation
\[
\ell_0^n \ell_1^{-n}=\varepsilon_0^n\sigma_\infty^{-n}\varepsilon_1^n\sigma_\infty^n,
\]
i.e. (comme $\ell_0=\varepsilon_0^2$, $\ell_1=\varepsilon_1^2$) si $n$ pair
\[
\varepsilon_0^n=\deleted{\unsure{\sigma_\infty}}\varepsilon_1^{3n}
\qquad\text{d'où }n=0,
\]
et si $n$ impair
\[
\varepsilon_0^n=\deleted{\unsure{\varepsilon_0}}\underbrace{\sigma_\infty\varepsilon_1^n\sigma_\infty}_{\sigma_\infty(\varepsilon_1)^n=\varepsilon_0^n}\varepsilon_1^{2n}
\]
i.e. $\varepsilon_1^{2n}=0$, i.e. encore $n=0$, cqfd.

\margin{\uline{NB} $\unsure{\mathcal M_{0,3}^{!\approx}(0)}$ \unsure{une} \illegible{2} $\mathcal N_{0,3}^{!\sim}$ tout entier, \unsure{au moins il} reste à voir si \unsure{l'autre}\ldots\ \illegible{1}}

\uline{\unsure{Scholie}} On a donc trouvé une section partielle canonique de $\mathcal M_{0,3}^{!\sim}$ au dessus du sous-groupe $\mathcal N_{0,3}^{!\approx}$ de $\mathcal N_{0,3}^{!\sim}$. Comme
\margin{236}

\scanpage{238}

ce sous-groupe contient $\Gamma_{\mathbb Q}$, on trouve donc une section de l'extension $\mathcal M_{0,3}^{!}$ de $\Gamma_{\mathbb Q}$ par $\widehat\Pi_{0,3}$, \deleted{Il faut} \inserted{qui est au fait une} section du sous-groupe de $\mathcal M_{0,3}^{!}$ normalisateur de $L_0$. Il faudrait vérifier que cette section canonique est associée au choix d'une uniformisante \unsure{sur le} point $0$ de $\widehat U_{0,3\mathbb Q}$, \deleted{\unsure{et qui sera déterminé}} \inserted{dont l'image dans} $\unsure{\mathfrak m_0}/\mathfrak m_0^2$ soit \inserted{bien} déterminée, où $\mathfrak m_0$ désigne l'idéal maximal de l'anneau local de $0$ dans $\widehat U_{0,3\mathbb Q}=\mathbb P_{\mathbb Q}^1$ \unsure{)} --- il faudrait déterminer cet élément de $\mathfrak m_0/\mathfrak m_0^2$ --- On a envie de prendre un élément qui soit, non seulement une base sur $\mathbb Q$, mais $\overline{\mathrm m}$ (dans la situation sur les entiers) une base sur $\mathbb Z$ --- \deleted{\unsure{ce}} \inserted{condition} qui détermine en effet un élément de $\mathfrak m_0/\mathfrak m_0^2$ \uuline{au} \uline{signe près} --- \unsure{ce ser.} en fait le choix correspondant à l'uniformisante $z$, ou $-z$. Mais est-ce bien le bon ? Il faudrait y revenir !
\margin{237}

\scanpage{239}
\mspage{420}

Travaillons désormais dans le sous-groupe $\mathcal M_{0,3}^{!\approx}$ de $\mathcal M_{0,3}^{!\sim}$, formé des $u$ qui commutent extérieurement à $\widehat\varphi$, i.e. image inverse de $\mathcal N_{0,3}^{!\approx}$, de sorte qu'on a
\[
\diagram{\begin{tikzcd}[ampersand replacement=\&]
1 \arrow[r] \& \widehat\Pi_{0,3} \arrow[r] \& \mathcal M_{0,3}^{!\approx} \arrow[r] \& \underset{\substack{\cup\\\Gamma_{\mathbb Q}}}{\mathcal N_{0,3}^{!\approx}} \arrow[r] \& 1
\end{tikzcd}}
\tag{24}
\]
On a donc trouvé un sous-groupe section
\[
\mathcal M_{0,3}^{!\approx}(0)\xrightarrow{\sim}\mathcal N_{0,3}^{!\approx}
\tag{25}
\]
\deleted{$\{\unsure{\xi}\mid\xi\in\widehat\Pi_{0,3}\text{ sol. de }(14)(15)\}$ \illegible{2}}
\[
\{u_\xi\mid\xi\in\widehat\Pi_{0,3}\text{ sol. de }(14)\text{ et }(15)\}\subset\mathcal M_{0,3}^{!\approx}.
\]

En termes de celui-ci on peut donc écrire $\mathcal M_{0,3}^{!\approx}$ comme produit $\frac12$ direct
\[
\mathcal M_{0,3}^{!\approx}\simeq\mathcal M_{0,3}^{!\approx}(\unsure{0})\cdot\widehat\Pi_{0,3},
\tag{26}
\]
en écrivant tout élément de $\mathcal M_{0,3}^{!\approx}$ de façon unique comme
\[
\begin{gathered}
\deleted{\unsure{u_{\xi,f}}}\operatorname{int}(f^{-1})\circ\unsure{u_0|u_\xi}\\
\bigl(\unsure{u_0|u_\xi}\in\mathcal M_{0,3}^{!\approx}(0)\text{ i.e. }u_0=u_\xi\text{ avec }\xi\text{ sol. de }(14)(15)\bigr)\\
f\in\widehat\Pi_{0,3}\text{ bien déterminé}
\end{gathered}
\tag{27}
\]
Ce sont justement les \deleted{$u_{\xi,f}=u_{\alpha,\beta,f}$} \inserted{$u_{\xi,f}$ de (19), soit,}
\[
u_{\xi,f}=u_{\alpha,\beta,f}
\tag{28}
\]
avec $\alpha,\beta$ déterminés en termes de $\xi$ par
\[
\alpha\rho(\alpha)^{-1}=\xi,\quad\deleted{\unsure{\varepsilon}}\beta\sigma_\infty(\beta)^{-1}=\xi \ell_1^{-\nu},
\tag{29}
\]
\margin{238}

\scanpage{240}

c'est ainsi qu'on peut interpréter les \unsure{triples} $(\alpha,\beta,f)$ \uline{normalisés} du cor. 1.

\uline{\unsure{Rescholie}} \inserted{(Récapitulation)} Le groupe $\mathcal N_{0,3}^{!\approx}\simeq\mathcal M_{0,3}^{!\approx}(0)$ est en corr. $1$-$1$ avec \unsure{l'ens.} des $\xi\in\widehat\Pi_{0,3}$ satisfaisant les deux relations
\[
\left\{
\begin{aligned}
\xi\rho(\xi)\rho^2(\xi)&=1\\
\xi&=\bigl(\widehat\varphi(\xi)\sigma_\infty(\widehat\varphi(\xi))^{-1}\bigr)\ell_1^\nu
\end{aligned}
\right.
\qquad\begin{gathered}\text{pour}\\\nu\in\widehat{\mathbb Z}\text{ conv.}\end{gathered}
\tag{30}
\]
ou encore (posant $\xi=\alpha\rho(\alpha)^{-1}$, $\alpha\in\unsure{\widehat{\widetilde\Pi}_{0,3}}$) avec \unsure{l'ens.} des $\alpha\in\unsure{\widehat{\widetilde\Pi}_{0,3}}$ satisfaisant la \inserted{deuxième} relation (30), dans laquelle on remplace $\xi$ par $\alpha\rho(\alpha)^{-1}$. L'élément $u_\xi$ de $\mathcal M_{0,3}^{!\approx}(0)$ défini par $\xi$ est défini par les formules (18), i.e. par
\[
u_\xi(\sigma_\infty)=(\xi \ell_1^{-\nu})\sigma_\infty,
\qquad u_\xi(\rho)=\xi\rho.
\tag{31}
\]
Enfin, si on a $u_\xi$ et $u_{\xi'}$, alors
\[
u_{\xi'}u_\xi=u_{\xi''},
\tag{32}
\]
où $\xi''$ est donné grâce aux formules p. \unsure{395'} (52) --- en notant qu'ici $\xi=g$, $\xi'=g'$, $\xi''=g''$ :
\deleted{$\xi''=\alpha''\rho(\alpha'')^{-1}$ avec $\alpha''=u_{\xi'}(\alpha)\cdot\alpha'$, $\unsure{f''=1}$}
\[
\deleted{\xi''=\deleted{\unsure{\alpha''\rho}}u_{\xi'}(\alpha)\underbrace{\alpha'\rho(\alpha')^{-1}}_{\xi'}\rho(u_{\xi'}(\alpha))^{-1}}
\]
\margin{239}

\scanpage{241}
\mspage{421}

\deleted{\illegible{1}}
\[
\xi''=u'(\xi)\cdot\xi'.
\tag{33}
\]

Posons maintenant
\[
\mathcal N_{0,3}^{\approx}
=\mathcal G_3\times\mathcal N_{0,3}^{!\approx}
\subset\mathcal N_{0,3}^{\sim}
=\mathcal G_3\times\mathcal N_{0,3}^{!\sim}
\tag{34}
\]
et soit $\mathcal M_{0,3}^{\approx}\subset\deleted{\unsure{\mathcal G}}\mathcal M_{0,3}^{\sim}$ l'image inverse dans $\mathcal M_{0,3}^{\sim}$ de $\mathcal N_{0,3}^{\approx}\subset\mathcal N_{0,3}^{\sim}$, de sorte qu'on a une structure d'extension
\[
1\longrightarrow\deleted{\unsure{\widehat{\mathcal G}_{0,3}^{+}}}\widehat\Pi_{0,3}
\longrightarrow\mathcal M_{0,3}^{\approx}
\unsure{\twoheadrightarrow|\longrightarrow}
\underset{\simeq\mathcal G_3\times\mathcal N_{0,3}^{!\approx}}{\mathcal N_{0,3}^{\approx}}
\longrightarrow1
\tag{35}
\]
contenant l'\deleted{extension}\inserted{\unsure{sous-extension}} (24) comme l'image inverse de \deleted{\unsure{la}}
\deleted{\illegible{7}}
\[
\mathcal N_{0,3}^{!\approx}
=\operatorname{Ker}(\mathcal N_{0,3}^{\approx}\longrightarrow\mathcal G_3).
\]
\deleted{\unsure{(36)}} \inserted{\unsure{On a}} \deleted{\unsure{aussi}} une \unsure{autre} \unsure{str. d'ext. sur} $\mathcal M_{0,3}^{\approx}$
\[
1\longrightarrow\widehat{\mathcal G}_{0,3}^{+}
\longrightarrow\mathcal M_{0,3}^{\approx}
\longrightarrow\mathcal N_{0,3}^{!\approx}\longrightarrow1,
\tag{36}
\]
contenant (24) comme sous-extension de $\overline{\mathrm m}$ quotient par un groupe plus petit.

Le sous-groupe $\mathcal M_{0,3}^{!\approx}(0)$, qui est une section de (24), est donc aussi une section de (36), d'où un iso. can.
\[
\mathcal M_{0,3}^{\approx}
\simeq\underbrace{\unsure{\mathcal N|\mathcal M}_{0,3}^{!\approx}(0)}_{\simeq\mathcal N_{0,3}^{!\approx}}
\cdot\widehat{\mathcal G}_{0,3}^{+}
\tag{37}
\]
(produit $\frac12$-direct).
\margin{240}

\scanpage{242}

Considérons maintenant l'application
\[
\Phi:\quad
\diagram{\begin{tikzcd}[ampersand replacement=\&,column sep=large]
\mathcal M_{0,3}^{!\approx}\arrow[r]\arrow[d,equal]
\&\mathcal M_{0,3}^{\approx}\arrow[d,equal]\\
\unsure{\mathcal M}_{0,3}^{!\approx}(0)\cdot\widehat\Pi_{0,3}
\&\mathcal M_{0,3}^{!\approx}(0)\cdot\widehat{\mathcal G}_{0,3}^{+}
\end{tikzcd}}
\tag{38}
\]
qui est \unsure{donnée|donné} par
\[
\Phi(u\cdot a)=u\widehat\varphi(a)
\tag{39}
\]
($u\in\mathcal M_{0,3}^{!\approx}(0)$, $a\in\widehat\Pi_{0,3}$). Il résulte de la caractérisation de $\mathcal M_{0,3}^{!\approx}(0)$ (cor. 3) --- \deleted{\unsure{i.e.}} \inserted{plus \uline{particulièrement}} (du fait que les éléments de $\mathcal M_{0,3}^{!\approx}(0)$ commutent effectivement (pas seulement extérieurement) à $\widehat\varphi$, que $\Phi$ est un homomorphisme de groupes. Ce n'est autre bien sûr que l'hom. $\Phi$ envisagé dans le \S\,41 (p.\,316') \inserted{(à restriction près aux sous-groupes $\mathcal M_{0,3}^{!}=\unsure{\mathcal E}_{0,3}$ et $\mathcal M_{0,3}$)} --- dans loc. cit. cet hom. pouvait se définir d'ailleurs de façon effective --- pas seulement à $\widehat\Pi_{0,3}$-conj. près ou à $\widehat{\mathcal G}_{0,3}^{+}$-conj. près --- et avec un peu de soin on devrait pouvoir vérifier \uline{l'égalité} de ces homomorphismes $\Phi$ du \S\,41 et d'ici, de façon \uline{effective}.
\margin{241}

\scanpage{243}
\mspage{422}
\uline{Calcul de} $x=\ell_\infty$.

Je vais revenir sur la détermination des $u_\xi\in\mathcal M_{0,3}^{!\sim}[0]=\operatorname{Norm}_{\mathcal M_{0,3}^{!\sim}}(L_0)$ qui sont \uline{normalisés} i.e. qui commutent à $\widehat\varphi$ i.e. satisfont (5) avec $\lambda=1$. Pour $x=\ell_0$ cette formule, en vertu de la forme particulière de $u_\xi$, est tautologique
\[
(\widehat\varphi u\ell_0=\widehat\varphi(\ell_0^\mu)=\varepsilon_0^\mu,
\quad u\widehat\varphi(\ell_0)=u(\varepsilon_0)=\varepsilon_0^\mu),
\]
et on \unsure{veut|voit} qu'il \unsure{revienne|revient} au $\overline{\mathrm m}$ de la poser pour $x=\ell_1$, ou pour $x=\ell_\infty$. Pour $x=\ell_1$, la condition prend la forme (10) --- remarquable du fait qu'elle donne un calcul \unsure{du|de} \deleted{la \unsure{variable}}\inserted{\unsure{paramètre}} ``délicate'' $\beta$ (dont l'existence \unsure{au priori|a priori}, en termes de $\xi$, repose sur des fondements qui n'ont pas encore été établis\ldots), en termes du paramètre ``grossier'' $\xi$ --- compte tenu du fait que $i$\inserted{$\in\mathbb Z/2\mathbb Z$} lui-$\overline{\mathrm m}$ est déterminé aussi aisément en termes de $\xi$.

Nous allons donner à cette condition une forme en apparence assez différente, en écrivant la relation de commutation (5) pour $x=\ell_\infty$, au lieu de $x=\ell_1$. Notons d'abord qu'on a la relation
\[
\tau'_\infty(\rho)=\rho^{-1}
\tag{40}
\]
Car
\[
\tau'_\infty(\rho)
=\sigma_\infty(\underbrace{\tau_\infty(\rho)}_{\ell_1^{-1}\rho})
=\deleted{\underbrace{\unsure{(\sigma_\infty\ell_1^{-1}\sigma_\infty^{-1})}}_{\unsure{\ell_0^{-1}}}}\inserted{\unsure{\ell_0^{-1}}}
\underbrace{\sigma_\infty(\rho)}_{\underbrace{\sigma_\infty\rho}_{\varepsilon_0}\sigma_\infty^{-1}}
=\varepsilon_0^{-1}\sigma_\infty^{-1}=\rho^{-1},
\]
\margin{242}

\scanpage{244}

il en résulte que si $u\in\mathcal M_{0,3}^{!\sim}$ satisfait
\[
u(\rho)=\rho^{\unsure{\mu|j}}
\qquad \unsure{\mu|j}\in\mathbb Z/3\mathbb Z\text{ convenable}
\tag{41}
\]
alors on a nécessairement $\unsure{j|\mu}\equiv-1\text{ ou }+1\ (3)$ et si \deleted{\unsure{$\alpha$ est} \illegible{1}$(j)=\alpha_3$ est l'entier \illegible{1}}
\[
\alpha_3(\unsure{j|\mu})\in\mathbb Z/2\mathbb Z
\quad\text{caractérisé par}\quad
\unsure{j|\mu}\equiv(-1)^{\alpha_3(\unsure{j|\mu})}\ (3)
\tag{42}
\]
on aura donc $u(\rho)=\rho$ ou $u(\rho)=\rho^{-1}$ suivant que $\alpha_3(\unsure{j|\mu})=0$ ou $\alpha_3(\unsure{j|\mu})=1$, et en tous cas on aura
\[
u(\rho)=\tau'_\infty{}^{\alpha_3}(\rho).
\tag{43}
\]

Ceci posé, faisons $x=\ell_\infty$ dans (5):
\[
\widehat\varphi u(\ell_\infty)
=\widehat\varphi(\operatorname{int}(f_\infty^{-1})\ell_\infty^\mu)
=\operatorname{int}(\widehat\varphi(f_\infty^{-1}))\rho^{-\mu}
\]
(car $\widehat\varphi(\ell_\infty)=\rho^{-1}$), et
\[
u(\widehat\varphi(\ell_\infty))
=u(\rho^{-1})=u(\rho)^{-1}
=\deleted{\unsure{(\xi\rho)^{-1}}=\rho^{-1}\xi^{-1}
=\unsure{(\rho^{-1}\xi^{-1}\rho)\rho^{-1}}}
\inserted{\operatorname{int}(\alpha)(\rho^{-1})}
\]
donc la formule (5) pour $\lambda=1$ devient
\[
\operatorname{int}(\widehat\varphi(f_\infty)\alpha)(\rho^{-1})
=\rho^{-\mu}\qquad\text{soit}
\]
\[
\operatorname{int}(\widehat\varphi(f_\infty)\alpha)\rho=\rho^\mu
\]
et en vertu de (43) \deleted{\unsure{ceci équiv.} $\unsure{\rho^\mu=\tau'_\infty{}^{\alpha_3}(\rho)}$}\inserted{\unsure{on a}} \deleted{\unsure{que}}
\[
\rho^\mu=\tau'_\infty{}^{\alpha_3=\alpha_3(\mu)}(\rho)
\overset{\unsure{\mathrm{df}|\mathrm{déf}}}{=}
\operatorname{int}(\tau'_\infty{}^{\alpha_3})(\rho)
\]
donc la condition devient
\[
\operatorname{int}(\tau'_\infty{}^{\alpha_3}\widehat\varphi(f_\infty)\alpha)\rho=\rho
\tag{\unsure{44}}
\]
i.e. que l'élément $\tau'_\infty{}^{\alpha_3}\widehat\varphi(f_\infty)\alpha$ du $\widehat{\mathcal G}_{0,3}=\widehat{\mathcal G}_{0,3}^{+}\cdot\{1,\tau_\infty\}$ commute à $\rho$.
\margin{243}
